\aircrafttype{de Havilland Vampire}

The de Havilland Vampire is a first-generation British jet day fighter, night fighter, and fighter-bomber. It first flew in 1943 but did not enter service until after the end of World War II. 

It features a single engine, straight wings, and an innovative twin-boom tail. The forward fuselage is fabricated from molded plywood, spruce, and balsa, building on de Havilland's extensive experience with this technique most notably in the Mosquito, but the wings, rear fuselage, and tail are aluminum for strength. Its armament is four Hispano 20 mm cannon, the standard for British fighters at that time, mounted under the nose. Initially, pilot was not provided with an ejector seat, although one was refitted to some export versions.

The Vampire entered service with the RAF in 1946. Its contemporary, the twin-engined Meteor, was more complex, expensive, and had less endurance on internal fuel. However, the Meteor crucially had a higher rate of climb, and this lead to it being selected for the interceptor role, and the Vampire served as a day fighter and fighter-bomber.

Nevertheless, the Vampire had considerable success on the export market, perhaps because of its lower price and simplicity compared to the Meteor, and for many air forces was their first jet aircraft.

The early F.1 and F.3 were day fighters. The F.3 had the slightly more powerful Goblin 2 engine (although this was also fitted to later F.1s) and significantly more internal fuel. Both could carry 100 imperial gallon (450L) fuel tanks under the outer wings, but had no provision for air-to-ground weapons other than their guns.

The FB.5 is the first fighter-bomber version. It was developed from the F.3 and featured clipped wings for better performance at low level, armor to protect the engine, provision for bombs in place of the fuel tanks, and rails under the inner wings for up to eight RP-3 rockets. The FB.6 and FB.9 were both derived from the FB.5, with the FB.6 featuring the more powerful Goblin 3 engine and the FB.9 having air conditioning for use in tropical zones (although Rhodesian FB.9s also were upgraded with the Goblin 3 for better performance). The FB.50 and FB.52 were licensed or export versions of the FB.6, and the FB.52A was a version of the FB.52 with the Goblin 2 engine. The Sea Vampire F.20 was adapted from the FB.5 for carrier-operation, and had a strengthened undercarriage and an arrestor hook.

The Vampire NF.10 is a night fighter version, derived from the FB.5 but with a new forward fuselage for the AI Mk X (SCR-720B) radar and its operator and the Goblin 3 engine to compensate for the weight of the radar equipment. It was developed initially for export, but was taken up by the RAF as to bridge the gap between the Mosquito NF.36 and the Meteor NF.11. The NF.54 was the corresponding export version.


The Sea Vampire F.20 is an adaptation of the FB.5 for RN carrier operations. It has strengthened undercarriage, an arrestor hook, and more effective speed brakes, but no capacity to fold its wings.

The Vampire T.11 is a trainer version of the FB.5, but retains full combat capability. RAF aircraft we refitted with ejection seats between 1954 and 1957. The Sea Vampire T.22 trainer is essentially a T.11 adapted to RN standards, but is not carrier-capable. It was flown exclusively from terrestrial air stations. RN aircraft were refitted with ejection seats in 1956 and 1957. The T.55 is the export version of the T.11.

The F.1 served exclusively and in small numbers with the RAF as a day fighter. The F.3 served more widely, with the RAF, Canadian RCAF, and Norwegian Luftforsvaret. The FB.5 was used by the RAF, Canadian RCAF, New Zealand RNZAF, South African SAAF, and Venezuelan FAV. The FB.6 served with the Swiss Flugwaffe. The FB.9 served with the RAF, Australian RAAF, New Zealand RNZAF, and Rhodesian SRAF/RRAF. The NF.10 served with the RAF. The FB.50 served with the Swedish Flygvapnet, the FB.52 with air forces of Ceylon, Egypt, Finland, India, Iraq, Italy, Jordan, Lebanon, New Zealand, Norway, Southern Rhodesia, Saudi Arabia, South Africa, Syria, Venezuela, and the FB.52A with the Italian AMI, the Syrian air force, and the Egyptian air force. 

The NF.10 served with the RAF, and the NF.54 served with the Italian AMI and the Indian IAF.

The T.11 server with the RAF. The T.55 served with the Indian IAF and the Austrian Luftstreitkräfte.

The Sea Vampire F.20 served with the RN, but for trials and familiarization and not as a front-line fighter. The T.22 also served as an advanced trainer.

The gun armament is four Hispano 20 mm cannons with 150 rounds per gun. A typical air-to-air load would be two fuel tanks to increase endurance. For air-to-ground missions, a typical load would be eight RP-3 rockets and then either two 500 lb bombs or fuel tanks, depending on the mission radius. On short-range missions, two 1,000-lb bombs could be carried but without rockets.

ADCs are provided for the:
\begin{adclist}
    \adcitem{Vampire F.3}
    \adcitem{Vampire FB.5}
    \adcitem{Vampire FB.6}
    \adcitem{Vampire FB.9}
    \adcitem{Vampire FB.9 (Goblin 3)}
    \adcitem{Vampire NF.10}
    \adcitem{Vampire T.11}
    \adcitem{Vampire T.11 (Late)}
    \adcitem{Sea Vampire F.20}
    \adcitem{Sea Vampire T.22}
    \adcitem{Sea Vampire T.22 (Late)}
    \adcitem{Vampire FB.50}
    \adcitem{Vampire FB.52}
    \adcitem{Vampire FB.52A}
    \adcitem{Vampire NF.54}
    \adcitem{Vampire T.55}
    \adcitem{Vampire T.55 (Late)}
\end{adclist}


See also:
\begin{typelist}
    \typeitem{DHA Vampire}
    \typeitem{FIAT Vampire}
    \typeitem{HAL Vampire}
    \typeitem{Macchi Vampire}
    \typeitem{SNCASE Vampire}
    \typeitem{SNCASE Mistral}
\end{typelist}